\documentclass{amsart}

% 数学相关宏包
\usepackage{amsmath,mathtools,amsthm,amssymb}
\usepackage{bbm}
\usepackage{mathrsfs}
\usepackage{centernot}

% 图形和绘图
\usepackage{graphicx,float}
\usepackage{tikz,tikz-cd}
\usepackage{pgfplots}
\usepackage[framemethod=TikZ]{mdframed}
\usetikzlibrary{decorations.pathmorphing,positioning,arrows,automata,calc,shapes,patterns}
\pgfplotsset{compat=1.18}
\usepgfplotslibrary{statistics}

% 表格和列表
\usepackage{booktabs,multirow,colortbl}
\usepackage{enumitem}
\setlist[itemize,enumerate,description]{noitemsep}
\setlistdepth{9}

% 算法
\usepackage{algorithm,algorithmic}

% 字体和颜色
\usepackage{xcolor,listings}
\usepackage{xeCJK}
\setCJKmainfont{Noto Serif CJK SC}

% 页面和链接
\usepackage{fancyhdr,caption}
\usepackage{hyperref}
\usepackage{cleveref}

% 工具宏包
\usepackage{etoolbox,letltxmacro,docmute}
\usepackage{gensymb}

% 自定义设置
\renewcommand{\arraystretch}{1.5}
\let\originalmathbb\mathbb
\renewcommand{\mathbb}[1]{%
    \ifstrequal{#1}{1}{\mathbbm{#1}}{\originalmathbb{#1}}%
}

% 颜色定义
\definecolor{lightblue}{RGB}{173,216,230}
\definecolor{lightgreen}{RGB}{144,238,144}
\definecolor{lightyellow}{RGB}{255,255,224}
\definecolor{lightgray}{RGB}{211,211,211}

% 超链接设置
\hypersetup{
    colorlinks=true,
    linkcolor=red,
    filecolor=magenta,
    urlcolor=cyan,
    pdftitle={\@title},
    pdfauthor={\@author},
    pdfsubject={数学笔记},
    pdfkeywords={数学},
    bookmarksnumbered=true,
    bookmarksopen=true,
    bookmarksopenlevel=6,
    pdfstartview=FitH,
    pdfpagemode=UseOutlines,
    pdfborder={0 0 0}
}

% 定理环境
\theoremstyle{plain}
\newtheorem{theorem}{Theorem}[section]
\newtheorem{lemma}[theorem]{Lemma}
\newtheorem{corollary}[theorem]{Corollary}
\newtheorem{proposition}[theorem]{Proposition}

\theoremstyle{definition}
\newtheorem{definition}[theorem]{Definition}
\newtheorem{example}[theorem]{Example}
\newtheorem{exercise}[theorem]{Exercise}
\newtheorem{problem}[theorem]{Problem}

\theoremstyle{remark}
\newtheorem{remark}[theorem]{Remark}
\newtheorem{note}[theorem]{Note}
\newtheorem{observation}[theorem]{Observation}

% 解答环境
\newtheorem*{solution}{Solution}
\newtheorem*{proof*}{Proof}

% 图片路径
\graphicspath{{F:/iCloudDrive/iCloud~md~obsidian/2025-summer/figures/}}

\author{乐绎华, Yue Yihua}
\address{中山大学, Sun Yat-sen University}
\email{yueyh@mail2.sysu.edu.cn}
\date{\today}
\title{Mathematical Notes}

\begin{document}

\maketitle
\tableofcontents

Given
\[
Y\subset\text{Gr}(k,n)\coloneqq \{ V_k\subseteq \mathbb{C}^{n}:\dim V_k=k \}
\]
ask
\[
H^{p,q}(Y)=?
\]
$H^{p,q}$ 通常指的是 \textbf{霍奇上同调群} (Hodge cohomology group)。它是在复几何和霍奇理论中一个非常重要的概念，用于研究复流形的上同调。

让我们从头开始，逐步解释其定义：

\section{微分形式的定义}

\subsection{1. 微分形式和外微分}

在一个光滑流形 $M$ 上，我们有微分形式 (differential forms)。一个 $k$ 阶微分形式 $\omega$ 可以写成局部坐标下的表达式，例如：
\[
\omega = \sum_{|I|=k} f_I(x) dx_{i_1} \wedge \cdots \wedge dx_{i_k}
\]
外微分算子 $d$ 将一个 $k$ 阶微分形式映射到一个 $(k+1)$ 阶微分形式，并且满足 $d^2 = 0$。

\subsection{2. de Rham 上同调}

de Rham 上同调 $H^k_{dR}(M)$ 是由闭形式 (closed forms，即 $d\omega = 0$) 的空间除以恰当形式 (exact forms，即 $\omega = d\eta$) 的空间得到的商群：
\[
H^k_{dR}(M) = \frac{\{\omega \in \Omega^k(M) \mid d\omega = 0\}}{\{d\eta \mid \eta \in \Omega^{k-1}(M)\}}
\]
de Rham 定理告诉我们，这个上同调群与流形的拓扑结构紧密相关，它同构于带实数系数的奇异上同调群 $H^k(M, \mathbb{R})$。

\subsection{3. 复流形上的微分形式}

当流形 $X$ 是一个\textbf{复流形}时，它的切丛具有一个复结构。这使得我们可以将微分形式进一步分解。
局部坐标系中，我们可以将复坐标 $z_j = x_j + i y_j$ 及其共轭 $\bar{z}_j = x_j - i y_j$ 引入，并定义微分算子：
\[
\partial = \sum_{j=1}^n \frac{\partial}{\partial z_j} dz_j, \quad \bar{\partial} = \sum_{j=1}^n \frac{\partial}{\partial \bar{z}_j} d\bar{z}_j
\]
任何一个微分形式 $\omega$ 都可以被分解为\textbf{类型} $(p,q)$ 的分量 $\omega^{p,q}$。一个 $(p,q)$ 型的微分形式在局部坐标下可以写成：
\[
\omega^{p,q} = \sum_{|I|=p, |J|=q} f_{I,J}(z, \bar{z}) dz_{i_1} \wedge \cdots \wedge dz_{i_p} \wedge d\bar{z}_{j_1} \wedge \cdots \wedge d\bar{z}_{j_q}
\]
其中 $p$ 是 $dz$ 的个数，$q$ 是 $d\bar{z}$ 的个数。
外微分 $d$ 也可以分解为两个算子：
\[
d = \partial + \bar{\partial}
\]
其中 $\partial$ 将 $(p,q)$ 型形式映射到 $(p+1,q)$ 型形式，而 $\bar{\partial}$ 将 $(p,q)$ 型形式映射到 $(p,q+1)$ 型形式。
由于 $d^2 = 0$，我们有 $\partial^2 = 0$, $\bar{\partial}^2 = 0$, 以及 $\partial\bar{\partial} + \bar{\partial}\partial = 0$。

\subsection{4. 霍奇上同调群 \texorpdfstring{$H^{p,q}$}{H^p,q} 的定义}

在复流形 $X$ 上，\textbf{霍奇上同调群 $H^{p,q}(X)$} 的定义为：
\[
H^{p,q}(X) = \frac{\ker(\bar{\partial}: \Omega^{p,q}(X) \to \Omega^{p,q+1}(X))}{\mathrm{Im}(\bar{\partial}: \Omega^{p,q-1}(X) \to \Omega^{p,q}(X))}
\]
这里 $\Omega^{p,q}(X)$ 表示 $X$ 上所有 $(p,q)$ 型光滑微分形式的空间。
这个定义是基于\textbf{Dolbeault 复形} (Dolbeault complex) 的，即：
\[
\cdots \xrightarrow{\bar{\partial}} \Omega^{p,q-1}(X) \xrightarrow{\bar{\partial}} \Omega^{p,q}(X) \xrightarrow{\bar{\partial}} \Omega^{p,q+1}(X) \xrightarrow{\bar{\partial}} \cdots
\]
$H^{p,q}(X)$ 是这个复形在 $\Omega^{p,q}(X)$ 位置的上同调。

\subsection{5. 霍奇分解定理 (Hodge Decomposition Theorem)}

对于一个\textbf{紧致凯勒流形} (compact Kähler manifold) $X$，霍奇理论给出了一个深刻的结论，即所谓的\textbf{霍奇分解定理}。它将 de Rham 上同调群与霍奇上同调群联系起来：
\[
H^k_{dR}(X) \otimes_{\mathbb{R}} \mathbb{C} \cong H^k(X, \mathbb{C}) \cong \bigoplus_{p+q=k} H^{p,q}(X)
\]
这个分解告诉我们，带复系数的 de Rham 上同调群 $H^k(X, \mathbb{C})$ 可以分解成霍奇上同调群的直和。

此外，霍奇分解还满足\textbf{霍奇对称性} (Hodge symmetry)：
\[
H^{p,q}(X) \cong H^{q,p}(X)
\]
这个同构由复共轭诱导。

\textbf{总结}：$H^{p,q}$ 是一个复流形上的上同调群，它捕捉了微分形式的类型信息。对于紧致凯勒流形，它们共同构成了流形的整个 de Rham 上同调群，并且具有重要的霍奇对称性。它们在代数几何、微分几何和数学物理中有着广泛的应用。

\section{层上同调的定义}

你提到的“层上的定义”指的是\textbf{层上同调} (sheaf cohomology) 的视角。这是定义霍奇上同调群 $H^{p,q}(X)$ 的另一种等价方式，特别是在代数几何中非常常用。

\subsection{1. 全纯微分形式的层 (Sheaf of Holomorphic Differential Forms)}

在复流形 $X$ 上，我们可以定义一个层的序列。
首先，我们有\textbf{全纯函数层} $\mathcal{O}_X$。对于 $X$ 上的任意开集 $U$，$\mathcal{O}_X(U)$ 是 $U$ 上所有全纯函数构成的环。
然后，我们可以定义\textbf{全纯 $p$-形式的层} $\Omega^p_X$。对于开集 $U$，$\Omega^p_X(U)$ 是 $U$ 上所有全纯 $p$-形式的空间。一个全纯 $p$-形式 $\omega$ 局部可以写成：
\[
\omega = \sum_{|I|=p} f_I(z) dz_{i_1} \wedge \cdots \wedge dz_{i_p}
\]
其中系数 $f_I(z)$ 都是 $U$ 上的全纯函数。

注意，全纯 $p$-形式的层 $\Omega^p_X$ 是一个 $\mathcal{O}_X$-模层。

\subsection{2. Dolbeault 上同调的层上同调定义}

对于一个复流形 $X$，霍奇上同调群 $H^{p,q}(X)$ 可以通过层上同调来定义。
具体来说，\textbf{$H^{p,q}(X)$ 同构于由全纯 $p$-形式层 $\Omega^p_X$ 的第 $q$ 个层上同调群}：
\[
H^{p,q}(X) \cong H^q(X, \Omega^p_X)
\]

为了理解这个定义，我们需要知道层上同调是如何计算的：

\begin{itemize}
	\item \textbf{层上同调群 $H^q(X, \mathcal{F})$} 是一个由层 $\mathcal{F}$ 诱导的\textbf{层上同调复形}的上同调群。
	\item 通常使用\textbf{细层} (fine sheaf) 的分辨率来计算层上同调。在复流形上，\textbf{层 $\Omega^{p,q}$ ($(p,q)$ 型光滑微分形式的层) 是一个细层}。
	\item 我们可以使用\textbf{Dolbeault 复形}作为 $\Omega^p_X$ 的一个细层分辨率。Dolbeault 复形是这样的：
\[
0 \to \Omega^p_X \xrightarrow{\bar{\partial}} \Omega^{p,1} \xrightarrow{\bar{\partial}} \Omega^{p,2} \xrightarrow{\bar{\partial}} \cdots
\]
这里的 $\Omega^{p,q}$ 是**$(p,q)$ 型光滑微分形式的层**，其在开集 $U$ 上的截面就是我们之前讨论的 $\Omega^{p,q}(U)$。
\end{itemize}

根据层上同调的一般理论，一个层的第 $q$ 个层上同调群可以通过其细层分辨率来计算。利用 Dolbeault 复形作为分辨率，我们得到：
\[
H^q(X, \Omega^p_X) \cong \frac{\ker(\bar{\partial}: \Gamma(X, \Omega^{p,q}) \to \Gamma(X, \Omega^{p,q+1}))}{\mathrm{Im}(\bar{\partial}: \Gamma(X, \Omega^{p,q-1}) \to \Gamma(X, \Omega^{p,q}))}
\]
其中 $\Gamma(X, \mathcal{F})$ 表示层 $\mathcal{F}$ 在整个空间 $X$ 上的\textbf{全局截面}。
$\Gamma(X, \Omega^{p,q})$ 正是 $X$ 上所有 $(p,q)$ 型光滑微分形式的空间 $\Omega^{p,q}(X)$。

因此，这个层上同调的定义
\[
H^{p,q}(X) \cong H^q(X, \Omega^p_X)
\]正是我们之前基于微分形式的定义\[
H^{p,q}(X) = \frac{\ker(\bar{\partial}: \Omega^{p,q}(X) \to \Omega^{p,q+1}(X))}{\mathrm{Im}(\bar{\partial}: \Omega^{p,q-1}(X) \to \Omega^{p,q}(X))}
\]
的精确描述。

\subsection{总结}

这两种定义是等价的，它们分别从不同的角度阐述了同一个数学对象。

\begin{itemize}
	\item \textbf{微分形式的定义}更直接，通过对微分形式的操作（$\bar{\partial}$ 算子）来构造上同调群。
	\item \textbf{层上同调的定义}则更抽象和普适。它将霍奇上同调群置于一个更广阔的理论框架中，即层论。这个视角在代数几何中尤为重要，因为它允许我们将全纯函数和微分形式作为层来处理，并使用强大的层上同调工具来研究它们的性质。例如，\textbf{小平消灭定理 (Kodaira vanishing theorem)} 这样的重要结果就是用层上同调的语言表述和证明的。
\end{itemize}

\section{例子：\texorpdfstring{$\mathbb{P}^{1}(\mathbb{C})$}{P^1(C)}}

好的，我们来举一个具体的例子，计算最简单的紧致凯勒流形之一——\textbf{复射影直线} $\mathbb{P}^1(\mathbb{C})$ 的霍奇上同调群。

$\mathbb{P}^1(\mathbb{C})$ 也可以被看作是\textbf{黎曼球面} $S^2$。

\subsection{1. 流形的基本信息}

\begin{itemize}
	\item \textbf{复流形}: $\mathbb{P}^1$ 是一个复维度为 1 的流形。
	\item \textbf{凯勒流形}: 它是紧致凯勒流形。
	\item \textbf{复维度}: $n=1$。
	\item \textbf{非零上同调群}: 对于一个复维度为 $n$ 的流形 $X$，非零的霍奇上同调群 $H^{p,q}(X)$ 只有当 $0 \le p, q \le n$ 时才存在。对于 $\mathbb{P}^1$，这意味着我们只需要考虑 $p, q \in \{0, 1\}$。
\end{itemize}

因此，我们要计算的霍奇上同调群是：$H^{0,0}(\mathbb{P}^1)$, $H^{1,0}(\mathbb{P}^1)$, $H^{0,1}(\mathbb{P}^1)$, 和 $H^{1,1}(\mathbb{P}^1)$。

\subsection{2. 计算 \texorpdfstring{$H^{0,0}(\mathbb{P}^1)$}{H^0,0(P^1)}}

根据定义，$H^{0,0}(\mathbb{P}^1) \cong H^0(\mathbb{P}^1, \Omega^0_{\mathbb{P}^1}) = H^0(\mathbb{P}^1, \mathcal{O}_{\mathbb{P}^1})$。

\begin{itemize}
	\item $\mathcal{O}_{\mathbb{P}^1}$ 是 $\mathbb{P}^1$ 上的\textbf{全纯函数层}。
	\item $H^0(\mathbb{P}^1, \mathcal{O}_{\mathbb{P}^1})$ 的元素是 $\mathbb{P}^1$ 上的\textbf{全局全纯函数}。
	\item \textbf{刘维尔定理} (Liouville's Theorem) 在复分析中的一个推广指出，\textbf{任何定义在紧致复流形上的全纯函数都必须是常数}。
	\item 因此，$\mathbb{P}^1$ 上的全局全纯函数只有常数函数。
	\item 所以，$H^{0,0}(\mathbb{P}^1)$ 是一个 1 维复向量空间，其基由常数函数生成。
	\item \textbf{结论}: $\dim H^{0,0}(\mathbb{P}^1) = 1$。
\end{itemize}

\subsection{3. 计算 \texorpdfstring{$H^{1,0}(\mathbb{P}^1)$}{H^1,0(P^1)}}

根据定义，$H^{1,0}(\mathbb{P}^1) \cong H^0(\mathbb{P}^1, \Omega^1_{\mathbb{P}^1})$。

\begin{itemize}
	\item $\Omega^1_{\mathbb{P}^1}$ 是 $\mathbb{P}^1$ 上的\textbf{全纯 1-形式层}。
	\item $H^0(\mathbb{P}^1, \Omega^1_{\mathbb{P}^1})$ 的元素是 $\mathbb{P}^1$ 上的\textbf{全局全纯 1-形式}。
	\item 在局部坐标 $z$ 下，一个全纯 1-形式可以写成 $f(z) dz$。
	\item 我们可以用两个坐标图 $U_0$ 和 $U_1$ 来覆盖 $\mathbb{P}^1$。
	\begin{itemize}
		\item 在 $U_0$ 上使用坐标 $z$。
		\item 在 $U_1$ 上使用坐标 $w = 1/z$。
		\item 全纯 1-形式在交集 $U_0 \cap U_1$ 上的过渡条件是：
\[
dw = d(1/z) = -(1/z^2)dz = -w^2 dz
\]
	\end{itemize}
	\item 假设 $\omega = f(z)dz$ 是一个全局全纯 1-形式。在 $U_1$ 上，它必须表示为 $\omega = g(w)dw$，其中 $g(w)$ 也是全纯的。
	\item 代入过渡条件，我们得到 $f(z)dz = g(w)(-w^2 dz)$，所以 $f(z) = -w^2 g(w) = -(1/z^2)g(1/z)$。
	\item 为了让 $f(z)$ 在 $z=0$ 处全纯，$-(1/z^2)g(1/z)$ 必须在 $z=0$ 处没有奇点。
	\item 如果 $g(w)$ 是一个全纯函数，在 $w=\infty$ 处没有极点，那么 $g(w)$ 必须是常数。但如果 $g(w)=c \ne 0$，那么 $f(z)$ 将在 $z=0$ 处有一个 2 阶极点，这与 $f(z)$ 在 $U_0$ 上全纯矛盾。
	\item 唯一的可能性是 $f(z) = 0$。
	\item \textbf{结论}: 不存在非零的全局全纯 1-形式。因此，$\dim H^{1,0}(\mathbb{P}^1) = 0$。
\end{itemize}

\subsection{4. 计算 \texorpdfstring{$H^{0,1}(\mathbb{P}^1)$}{H^0,1(P^1)}}

\begin{itemize}
	\item 对于紧致凯勒流形，我们有\textbf{霍奇对称性}：$H^{p,q}(X) \cong H^{q,p}(X)$。
	\item 因此，$H^{0,1}(\mathbb{P}^1) \cong H^{1,0}(\mathbb{P}^1)$。
	\item \textbf{结论}: $\dim H^{0,1}(\mathbb{P}^1) = \dim H^{1,0}(\mathbb{P}^1) = 0$。
\end{itemize}

\subsection{5. 计算 \texorpdfstring{$H^{1,1}(\mathbb{P}^1)$}{H^1,1(P^1)}}

\begin{itemize}
	\item 这里我们使用霍奇分解定理，它将 de Rham 上同调与霍奇上同调联系起来：
\[
H^k_{dR}(\mathbb{P}^1, \mathbb{C}) \cong \bigoplus_{p+q=k} H^{p,q}(\mathbb{P}^1)
\]
	\item 我们知道 $\mathbb{P}^1 \cong S^2$ 的 de Rham 上同调群为：
	\begin{itemize}
		\item $\dim H^0_{dR}(\mathbb{P}^1, \mathbb{C}) = 1$
		\item $\dim H^1_{dR}(\mathbb{P}^1, \mathbb{C}) = 0$
		\item $\dim H^2_{dR}(\mathbb{P}^1, \mathbb{C}) = 1$
		\item 所有其他维度的都为 0。
	\end{itemize}
	\item 对于 $k=2$，我们有 $H^2_{dR}(\mathbb{P}^1, \mathbb{C}) \cong H^{2,0}(\mathbb{P}^1) \oplus H^{1,1}(\mathbb{P}^1) \oplus H^{0,2}(\mathbb{P}^1)$。
	\item 但是，我们已经知道 $p,q \le 1$，所以 $H^{2,0}=0$ 和 $H^{0,2}=0$。
	\item 因此，$H^2_{dR}(\mathbb{P}^1, \mathbb{C}) \cong H^{1,1}(\mathbb{P}^1)$。
	\item \textbf{结论}: $\dim H^{1,1}(\mathbb{P}^1) = \dim H^2_{dR}(\mathbb{P}^1, \mathbb{C}) = 1$。
\end{itemize}

\subsection{总结霍奇菱形}

我们将计算结果总结成 $\mathbb{P}^1$ 的霍奇菱形：
\[
\begin{array}{ccc}
& 1 & \\
0 & & 0 \\
& 1 &
\end{array}
\]
这表示：

\begin{itemize}
	\item $h^{0,0} = 1$
	\item $h^{1,0} = 0, h^{0,1} = 0$
	\item $h^{1,1} = 1$
\end{itemize}

这个计算结果与我们从 $\mathbb{P}^n$ 的一般性质（$h^{p,p}=1$, 其他为0）中推导出的结果是吻合的。

\textbf{这个例子展示了如何通过结合以下方法来计算霍奇上同调群：}

\begin{enumerate}
	\item 利用\textbf{基本定理}（如刘维尔定理）和流形的拓扑性质来计算低阶上同调群。
	\item 利用\textbf{流形上的微分形式}的性质和坐标转换来证明某些群为零。
	\item 利用\textbf{霍奇分解定理}和\textbf{霍奇对称性}来简化高阶上同调群的计算。
\end{enumerate}

这些方法通常比直接解微分方程要有效得多。

\section{为什么只计算出维数，而不写出具体的群？}

这是一个非常好的、也很深刻的问题。你的观察非常准确。

我之前提供的例子，包括 $\mathbb{P}^1$ 和格拉斯曼流形，确实都只给出了霍奇上同调群的\textbf{维数} $h^{p,q}$，而不是具体的群本身。

原因在于，对于大多数数学家而言，霍奇上同调群的\textbf{维数}和\textbf{结构}（例如，作为向量空间的基）比“群本身”更重要，更具有物理和几何意义。下面我来详细解释为什么会这样：

\subsection{1. 霍奇上同调群的本质}

\begin{itemize}
	\item 霍奇上同调群 $H^{p,q}(X)$ 本质上是一个\textbf{复向量空间}。
	\item 我们计算的 $h^{p,q}(X)$ 是这个向量空间的维数。维数是一个不变量，它描述了空间的大小。
	\item 说“计算出群”本身，通常指的是找到它的具体结构，例如它的加法规则，或者找到它的一组生成元。对于向量空间来说，这就意味着找到一个基。
\end{itemize}

\subsection{2. 计算基的困难性}

要找到霍奇上同调群的一组基，就必须找到一组\textbf{具体的}、\textbf{线性无关的}、\textbf{$(p,q)$ 型的}、\textbf{$\bar{\partial}$-闭的}微分形式，并且它们不等于任何 $\bar{\partial}$-恰当形式。

\begin{itemize}
	\item \textbf{$\bar{\partial}$-闭形式}: 要找到一个 $\bar{\partial}$-闭形式 $\omega$，即 $\bar{\partial}\omega=0$，本质上是要求解一个偏微分方程。这通常是极具挑战性的。
	\item \textbf{商空间}: 即使我们找到了一个 $\bar{\partial}$-闭形式，它也只代表了霍奇上同调群中的一个\textbf{同伦类}（一个上同调类）。要确定它是不是一个非零的同伦类，我们必须证明它不能写成 $\omega = \bar{\partial}\eta$ 的形式。这又是另一个解微分方程的问题。
	\item \textbf{规范形式}: 霍奇定理提供了一个很好的解决方案：每个同伦类都包含一个唯一的\textbf{调和形式}。所以，如果我们要找到一个基，我们应该寻找一组线性无关的调和 $(p,q)$ 形式。但是，解出调和方程 $\Delta_{\bar{\partial}} \omega = 0$ 同样非常困难。
\end{itemize}

\subsection{3. 维数的几何和拓扑意义}

相比之下，维数 $h^{p,q}$ 本身就承载了大量的几何和拓扑信息，是研究流形的重要工具。

\begin{itemize}
	\item \textbf{霍奇分解}: $H^k(X, \mathbb{C}) \cong \bigoplus_{p+q=k} H^{p,q}(X)$。计算出所有的 $h^{p,q}$，我们就知道了流形的贝蒂数 $b_k = \sum_{p+q=k} h^{p,q}$，从而知道了流形的欧拉示性数等重要的拓扑不变量。
	\item \textbf{霍奇对称性}: $h^{p,q} = h^{q,p}$。这对于紧致凯勒流形是成立的，如果一个流形不满足这个对称性，我们就知道它不是凯勒流形。
	\item \textbf{几何意义}: 许多重要的几何不变量，例如\textbf{算术亏格} (arithmetic genus) $\chi(X)$，可以用霍奇数来表示：
\[
\chi(X) = \sum_{p,q} (-1)^{p+q} h^{p,q}(X)
\]
	\item \textbf{代数几何}: 在代数几何中，霍奇数可以与\textbf{代数簇}上的代数环、模空间等结构联系起来。例如，\textbf{陈类} (Chern class) 可以用霍奇数来计算。
\end{itemize}

\subsection{4. 特殊情况下的基}

在某些非常特殊和简单的情况下，我们确实可以找到霍奇上同调群的基：

\begin{itemize}
	\item \textbf{$H^{0,0}(X)$}: 如果 $X$ 是连通的，那么 $H^{0,0}(X)$ 的维数是 1，基就是常数函数 1。这是一个非常具体的基。
	\item \textbf{黎曼面 $X$}: $H^{1,0}(X)$ 的维数是亏格 $g$。我们可以找到由 $g$ 个全纯 1-形式构成的基，这些形式在某些正则坐标下有非常明确的表达式。这个计算通常在黎曼面理论中完成。
\end{itemize}

\subsection{总结}

\begin{itemize}
	\item 大部分情况下，我们计算霍奇上同调群的\textbf{维数} $h^{p,q}$，因为这是一个可以计算出的、有重要几何和拓扑意义的不变量。
	\item 维数本身已经足够回答大部分数学问题了。
	\item 具体找到一个基（即找到具体的微分形式）通常是一个非常困难的分析问题，远超出了大多数研究的需要，除非研究本身就是关于这些微分形式的性质。
	\item 你提出的问题非常敏锐，它触及了上同调理论中抽象代数结构（商群）与具体分析对象（微分形式）之间的核心区别。在实践中，数学家更关心上同调群作为一个抽象对象的性质，而维数是描述这些性质的最重要指标之一。
\end{itemize}

好的，我将根据你提出的要求，对之前的例子 $\mathbb{P}^1$ 给出具体的霍奇上同调群，而不仅仅是它们的维数。这将涉及到寻找上同调群的\textbf{生成元}，即作为向量空间的基。

\section{例子：\texorpdfstring{$\mathbb{P}^1(\mathbb{C})$}{P^1(C)} 的具体霍奇上同调群}

我们来回顾一下 $\mathbb{P}^1$ 的霍奇数：$h^{0,0}=1$, $h^{1,1}=1$，其余均为 0。

\subsection{1. \texorpdfstring{$H^{0,0}(\mathbb{P}^1)$}{H^0,0(P^1)}}

\begin{itemize}
	\item \textbf{维数}: $\dim H^{0,0}(\mathbb{P}^1) = 1$。
	\item \textbf{群}: $H^{0,0}(\mathbb{P}^1) = H^0(\mathbb{P}^1, \mathcal{O}_{\mathbb{P}^1})$ 是全局全纯函数构成的空间。
	\item \textbf{基}: 根据刘维尔定理，全局全纯函数只能是常数。所以，这个向量空间的一组基可以由\textbf{常数函数 $1$} 来表示。
	\item \textbf{具体的群}:
\[
H^{0,0}(\mathbb{P}^1) \cong \mathbb{C} \cdot 1
\]
这是一个一维的复向量空间，由函数 $f(z) = c$（其中 $c \in \mathbb{C}$）生成。
\end{itemize}

\subsection{2. \texorpdfstring{$H^{1,0}(\mathbb{P}^1)$}{H^1,0(P^1)} 和 \texorpdfstring{$H^{0,1}(\mathbb{P}^1)$}{H^0,1(P^1)}}

\begin{itemize}
	\item \textbf{维数}: $\dim H^{1,0}(\mathbb{P}^1) = 0$ 和 $\dim H^{0,1}(\mathbb{P}^1) = 0$。
	\item \textbf{群}: 零维向量空间，即\textbf{零空间} $\{0\}$。
	\item \textbf{具体的群}:
\[
H^{1,0}(\mathbb{P}^1) \cong \{0\}
\]
\[
H^{0,1}(\mathbb{P}^1) \cong \{0\}
\]
它们没有非零的生成元。
\end{itemize}

\subsection{3. \texorpdfstring{$H^{1,1}(\mathbb{P}^1)$}{H^1,1(P^1)}}

\begin{itemize}
	\item \textbf{维数}: $\dim H^{1,1}(\mathbb{P}^1) = 1$。
	\item \textbf{群}: $H^{1,1}(\mathbb{P}^1)$ 是由 $(1,1)$ 型的\textbf{调和形式}构成的空间。我们需要找到一个具体的、非零的调和 $(1,1)$ 形式。
	\item \textbf{基}: 一个经典的例子是 \textbf{Fubini-Study 度量} 的凯勒形式 $\omega_{FS}$。
	\begin{itemize}
		\item 我们可以使用坐标图 $U_0 = \{[z_0:z_1] \mid z_0 \neq 0\}$ 和 $U_1 = \{[z_0:z_1] \mid z_1 \neq 0\}$ 来覆盖 $\mathbb{P}^1$。
		\item 在 $U_0$ 上使用坐标 $z = z_1/z_0$，Fubini-Study 凯勒形式 $\omega_{FS}$ 可以写成：
\[
\omega_{FS} = \frac{i}{2\pi} \partial\bar{\partial} \log(1+|z|^2) = \frac{i}{2\pi} \frac{dz \wedge d\bar{z}}{(1+|z|^2)^2}
\]
		\item 这个形式 $\omega_{FS}$ 满足以下重要的性质：
		\begin{enumerate}
			\item 它是\textbf{全局定义}在 $\mathbb{P}^1$ 上的一个光滑 $(1,1)$ 形式。
			\item 它是\textbf{闭形式}，即 $d\omega_{FS}=0$。在凯勒流形上，这等价于 $\partial\omega_{FS}=0$ 和 $\bar{\partial}\omega_{FS}=0$。
			\item 它代表了一个非零的\textbf{de Rham 上同调类} $[\omega_{FS}] \in H^2_{dR}(\mathbb{P}^1, \mathbb{R})$。这个类是 $\mathbb{P}^1$ 的\textbf{基本类}，其积分 $\int_{\mathbb{P}^1} \omega_{FS} = 1$。
		\end{enumerate}
		\item 由于 $\omega_{FS}$ 是一个闭形式，它也是 $\bar{\partial}$-闭的。
		\item 因为 $\mathbb{P}^1$ 是紧致凯勒流形，霍奇定理告诉我们这个上同调类 $[\omega_{FS}]$ 对应于一个唯一的调和 $(1,1)$ 形式。事实上，$\omega_{FS}$ 本身就是调和的。
		\item 所以，$\omega_{FS}$ 构成了 $H^{1,1}(\mathbb{P}^1)$ 的一组基。
	\end{itemize}
	\item \textbf{具体的群}:
\[
H^{1,1}(\mathbb{P}^1) \cong \mathbb{C} \cdot [\omega_{FS}]
\]
其中 $[\omega_{FS}]$ 代表由 Fubini-Study 凯勒形式生成的上同调类。
\end{itemize}

\subsection{总结}

通过这个例子，我们可以看到“具体的群”是什么意思。

\begin{itemize}
	\item 对于 $H^{0,0}$，基就是具体的函数 $f(z) = 1$。
	\item 对于 $H^{1,1}$，基是具体的微分形式 $\omega_{FS}$ 所代表的\textbf{上同调类} $[\omega_{FS}]$。
\end{itemize}

值得注意的是，在大多数情况下，我们找到的基并不是一个单一的函数或形式，而是一个\textbf{由这个函数或形式代表的上同调类}。这是因为上同调群的定义本身就是基于商空间的。在 $H^{1,1}(\mathbb{P}^1)$ 的例子中，我们找到了一个具体的代表元 $\omega_{FS}$，它就是霍奇定理保证存在的那个唯一的调和形式。

这个例子很好地说明了在简单流形上，我们确实可以具体地写出这些上同调群的生成元，但这通常需要利用流形的一些特殊几何结构（例如，凯勒形式）。对于更复杂的流形，即使维数已知，找到具体的基形式仍然是一个巨大的挑战。

\section{Kähler Manifold}

\href{https://math.uchicago.edu/~dannyc/notes/kahler_notes.pdf}{kahler\_notes.pdf}

\href{https://www.mathi.uni-heidelberg.de/~freitag/skripten/kaehler.pdf}{kaehler.pdf}

Sure, building on our discussion of Grassmannians and their cohomology, I'll outline a clear, general process to calculate the dimension of the $(p,p)$-cohomology group $H^{p,p}(\mathrm{Gr}(k,n))$ for the complex Grassmannian $\mathrm{Gr}(k,n)$. This dimension is the Hodge number $h^{p,p}(\mathrm{Gr}(k,n))$, and since the Hodge structure is diagonal (as we discussed), it's also equal to the $2p$-th Betti number $b_{2p}(\mathrm{Gr}(k,n))$.

Recall that for $\mathrm{Gr}(k,n)$, all off-diagonal Hodge groups vanish ($H^{p,q} = 0$ for $p \neq q$), so we only need to focus on these diagonal terms. The process relies on Schubert calculus and partition counting, which is combinatorial and straightforward to compute for specific $k$, $n$, and $p$.

\subsection{General Step-by-Step Process}

\begin{enumerate}
	\item \textbf{Identify the Parameters and Bounds:}
	\begin{itemize}
		\item Specify $k$ (the dimension of subspaces) and $n$ (the ambient dimension), with $1 \leq k < n$.
		\item Note that the total complex dimension of $\mathrm{Gr}(k,n)$ is $k(n-k)$, so $p$ ranges from 0 to $k(n-k)$. For $p$ outside this range, $h^{p,p} = 0$.
		\item The Schubert classes are indexed by partitions (Young diagrams) that fit inside a rectangle of height $k$ and width $n-k$.
	\end{itemize}
	\item \textbf{Count Relevant Partitions (Direct Combinatorial Method):}
	\begin{itemize}
		\item A partition $\lambda$ is a Young diagram where the number of boxes is $|\lambda| = p$, and it must fit entirely within the $k \times (n-k)$ rectangle (i.e., no more than $k$ rows, each of length at most $n-k$).
		\item List or enumerate all such distinct partitions $\lambda$ with exactly $p$ boxes.
		\item The number of these partitions is exactly $h^{p,p}(\mathrm{Gr}(k,n))$:
\[
h^{p,p}(\mathrm{Gr}(k,n)) = \#\{\lambda \subseteq k \times (n-k) : |\lambda| = p\}.
\]		\item For small values, you can do this manually by drawing diagrams. For larger ones, use combinatorial algorithms or software (e.g., SageMath or Mathematica has built-in partition counters).
	\end{itemize}
	\item \textbf{Use the Generating Function (q-Binomial Coefficient Method):}
	\begin{itemize}
		\item Alternatively, compute the Gaussian binomial coefficient (q-binomial), which generates all $h^{p,p}$ as coefficients:
\[
\binom{n}{k}_t = \prod_{i=1}^{k} \frac{1 - t^{n-k+i}}{1 - t^i} = \sum_{p=0}^{k(n-k)} h^{p,p}(\mathrm{Gr}(k,n)) \, t^p.
\]		\item Expand this polynomial in $t$ (analytically or using a computer algebra system) and extract the coefficient of $t^p$, denoted $[t^p] \binom{n}{k}_t$.
		\item This is equivalent to the partition count in step 2 but is often easier for computation, especially for larger $k$ and $n$.
	\end{itemize}
	\item \textbf{Verify and Interpret:}
	\begin{itemize}
		\item Confirm that $h^{p,p} = \dim H^{p,p}(\mathrm{Gr}(k,n))$, and since the cohomology is concentrated in even degrees, this also gives $b_{2p} = h^{p,p}$.
		\item Note the symmetry: $h^{p,p}(\mathrm{Gr}(k,n)) = h^{k(n-k)-p, k(n-k)-p}(\mathrm{Gr}(k,n))$, due to Poincaré duality.
	\end{itemize}
\end{enumerate}

\subsection{Example Computations}

\begin{itemize}
	\item For $\mathrm{Gr}(1, n) = \mathbf{CP}^{n-1}$: The rectangle is $1 \times (n-1)$. There's exactly one partition for each $p$ from 0 to $n-1$ (a single row of length $p$), so $h^{p,p} = 1$.
	\item For $\mathrm{Gr}(2,4)$: Rectangle $2 \times 2$, maximum $p=4$. Partition counts by $p$: 1 (empty), 1 (one box), 2 (two boxes: row or column), 1 (three boxes), 1 (four boxes). Thus, $h^{p,p} = (1,1,2,1,1)$ for $p=0$ to 4.
	\item For $\mathrm{Gr}(2,5)$: Rectangle $2 \times 3$, maximum $p=6$. Using the q-binomial:
\[
\binom{5}{2}_t = \frac{(1-t^4)(1-t^5)}{(1-t)(1-t^2)} = 1 + t + 2t^2 + 2t^3 + 2t^4 + t^5 + t^6,
\]so $h^{3,3} = 2$, for instance.
\end{itemize}

This process is general and works for any $k, n, p$; if you provide specific values, I can compute an example step by step! If you'd like references to tools or further details on Schubert calculus, let me know.


\end{document}